\documentclass[12pt]{article}

% ---------- packages ----------
\usepackage[margin=1in]{geometry}
\usepackage{amsmath,amsthm,amssymb,mathtools}
\usepackage{bbm}
\usepackage{booktabs}
\usepackage{enumitem}
\usepackage{natbib}
\usepackage{hyperref}
\usepackage{setspace}
\usepackage{cleveref}
\usepackage{graphicx}
\usepackage{xcolor}
\usepackage{subcaption}

\onehalfspacing

\hypersetup{
    colorlinks=true,
    linkcolor=blue!60!black,
    citecolor=blue!60!black,
    urlcolor=blue!60!black
}

% ---------- title ----------
\title{%
  \textbf{Online Appendix}\\[6pt]
  \large SaoMNK: A Stochastic Actor-Oriented Generalization\\
  of the NK Fitness Landscape Model\\[12pt]
  \normalsize Formal Proofs, Computational Demonstration,\\
  and Response to Anticipated Concerns}

\author{Stephen Downing\\
  \small University of Missouri}
\date{\today}

\begin{document}

\maketitle
\thispagestyle{empty}

\vspace{1cm}

\begin{abstract}
\noindent
This online appendix provides three components supporting the main text's
claim that the Stochastic Actor-Oriented Model on NK Landscapes (SaoMNK)
formally generalizes Kauffman's (1993) NK fitness landscape framework.

\medskip\noindent
\textbf{Appendix A} presents three formal theorems with complete proofs:
the \emph{Reduction Theorem} (NK is a degenerate special case of SaoMNK),
the \emph{Generalization Theorem} (SaoMNK strictly extends NK along three
independent dimensions), and the \emph{Approximation Theorem} (any NK
payoff structure can be reproduced within SaoMNK to arbitrary precision).

\medskip\noindent
\textbf{Appendix B} provides computational verification: an R script that
(i) implements both NK and SaoMNK side by side, (ii) verifies exact
numerical equivalence under the reduction conditions ($M=1$,
$\boldsymbol{\theta}=\mathbf{0}$, $\beta\to\infty$), and (iii) demonstrates
that $M>1$ produces endogenous landscape co-evolution effects impossible
in single-actor NK.

\medskip\noindent
\textbf{Appendix C} addresses eight anticipated reviewer concerns, with
detailed mathematical and conceptual responses.

\medskip\noindent
Together, these materials establish that SaoMNK is a micro-founded,
empirically estimable generalization of NK that preserves all of NK's
theoretical insights while adding multi-actor dynamics, endogenous
landscape co-evolution, and tunable bounded rationality.
\end{abstract}

\newpage
\tableofcontents
\newpage

% ============================================================
%  APPENDIX A: Formal Proofs
% ============================================================

\part*{Appendix A: Formal Proofs}
\addcontentsline{toc}{part}{Appendix A: Formal Proofs}

\noindent\textit{The complete formal proof document is contained in the
companion file} \texttt{saomnk\_nk\_equivalence\_proof.tex}\textit{.
When compiled for the final manuscript, this content should be included
here via} \verb|% ---------------------------------------------------------------------------
% NUMBERING AND STATUS NOTE (2026-09-09). This file's Theorems 1-3 are Paper T's
% Theorems 1-3 (CD2026_PaperT_v4, Management Science). Its "Proposition 1/2/3"
% under Theorem 3 are Paper T's Theorems 3a/3b/3c. Lemma 1 (Payoff Equivalence)
% was FALSE as written until 2026-09-09: the utility carried a factor b_{id} in
% front of f_d that NK fitness does not have, and the proof verified the
% identity only at x = 1. Repaired here and in Paper T (stage 127) by dropping
% the factor; the counterexample and the repaired proof are machine-checked in
% Lean 4 (D:\CD2026\formal\saomnk, Saomnk/Reduction.lean).
% ---------------------------------------------------------------------------
\documentclass[12pt]{article}

% ---------- packages ----------
\usepackage[margin=1in]{geometry}
\usepackage{amsmath,amsthm,amssymb,mathtools}
\usepackage{bbm}
\usepackage{booktabs}
\usepackage{enumitem}
\usepackage{natbib}
\usepackage{hyperref}
\usepackage{setspace}
\usepackage{cleveref}

\onehalfspacing

% ---------- theorem environments ----------
\newtheorem{theorem}{Theorem}
\newtheorem{proposition}[theorem]{Proposition}
\newtheorem{lemma}[theorem]{Lemma}
\newtheorem{corollary}[theorem]{Corollary}
\newtheorem{definition}[theorem]{Definition}
\newtheorem{remark}[theorem]{Remark}
\newtheorem{example}[theorem]{Example}

% ---------- operators and shortcuts ----------
\DeclareMathOperator*{\argmax}{arg\,max}
\newcommand{\R}{\mathbb{R}}
\newcommand{\N}{\mathbb{N}}
\newcommand{\E}{\mathbb{E}}
\newcommand{\bone}{\mathbbm{1}}
\newcommand{\bx}{\mathbf{x}}
\newcommand{\bb}{\mathbf{b}}
\newcommand{\be}{\mathbf{e}}
\newcommand{\bB}{\mathbf{B}}
\newcommand{\bA}{\mathbf{A}}
\newcommand{\bE}{\mathbf{E}}
\newcommand{\btheta}{\boldsymbol{\theta}}

% ---------- title ----------
\title{%
  \textbf{Online Appendix: Formal Proofs of the Relationship\\
  Between the SaoMNK Model and the NK Landscape}\\[6pt]
  \large Reduction, Generalization, and Approximation Results}

\author{}
\date{}

\begin{document}

\maketitle

\begin{abstract}
\noindent
This online appendix provides formal mathematical proofs establishing the
relationship between the Stochastic Actor-Oriented Model on NK Landscapes
(SaoMNK) and the standard NK model from complexity theory
\citep{kauffman1993origins, levinthal1997adaptation}. We prove three main
results. First, the \emph{Reduction Theorem} shows that the standard NK model is
a degenerate special case of SaoMNK, recoverable by restricting the model to a
single actor with deterministic greedy search. Second, the
\emph{Generalization Theorem} demonstrates that SaoMNK strictly extends NK along
three independent dimensions: multi-actor endogenous interaction, endogenous
landscape co-evolution, and bounded rationality with tunable search precision.
Third, the \emph{Approximation Theorem} establishes, via three distinct proof
strategies, that any NK payoff structure can be reproduced to arbitrary precision
within the SaoMNK framework. Together, these results position SaoMNK as a
micro-founded, empirically estimable generalization of the NK model for
organizational and strategic analysis.
\end{abstract}

\bigskip

% ============================================================
\section{Notation and Preliminaries}
\label{sec:notation}
% ============================================================

We begin by collecting the notation used throughout. \Cref{tab:notation}
summarizes all symbols.

\begin{table}[h!]
\centering
\caption{Notation summary.}
\label{tab:notation}
\small
\begin{tabular}{@{}lp{10.5cm}@{}}
\toprule
\textbf{Symbol} & \textbf{Description} \\
\midrule
$N$           & Number of binary components (decisions, policy choices) \\
$K$           & Epistasis parameter in the NK model; each component depends on $K$ others \\
$M$           & Number of actors (agents) in SaoMNK \\
$\bx \in \{0,1\}^N$ & Binary decision string (state) in the NK model \\
$\bB \in \{0,1\}^{M \times N}$ & Bipartite adjacency matrix in SaoMNK; row $\bb_i$ is actor $i$'s configuration \\
$\bA \in \{0,1\}^{N \times N}$ & Epistasis (interaction) matrix in the NK model \\
$\bE \in \{0,1\}^{N \times N}$ & Epistasis (interaction) matrix in SaoMNK \\
$\mathcal{N}(d)$ & Epistatic neighborhood of component $d$: $\mathcal{N}(d) = \{j : A_{dj} = 1\}$, $|\mathcal{N}(d)| = K+1$ \\
$f_d : \{0,1\}^{|\mathcal{N}(d)|} \to [0,1]$ & Component payoff function drawn i.i.d.\ $\mathrm{Uniform}(0,1)$ \\
$W(\bx)$      & NK fitness: $W(\bx) = \frac{1}{N}\sum_{d=1}^{N} f_d(\bx_{\mathcal{N}(d)})$ \\
$\mathrm{PK}(N)$ & Power key: $\mathrm{PK}(N) = (2^{N-1}, 2^{N-2}, \ldots, 2^0)$ \\
$\beta$       & Logit precision parameter (inverse temperature) \\
$\btheta$     & Vector of SAOM effect parameters \\
$s_k(\bb_i, \bB)$ & SAOM statistic $k$ for actor $i$ given full state $\bB$ \\
$u_i(\bb_i; \bB_{-i})$ & Actor $i$'s utility in SaoMNK \\
$\odot$       & Element-wise (Hadamard) product \\
\bottomrule
\end{tabular}
\end{table}

% ============================================================
\subsection{The NK Model}
\label{sec:nk_model}
% ============================================================

We state the NK model following \citet{kauffman1993origins} and
\citet{levinthal1997adaptation}, with notation consolidated from
\citet{altenberg1997nk}.

\begin{definition}[NK Landscape]
\label{def:nk}
An \emph{NK landscape} is a tuple $(N, K, \bA, \{f_d\}_{d=1}^N)$ where:
\begin{enumerate}[label=(\roman*)]
  \item $N \in \N$ is the number of binary components.
  \item $K \in \{0, 1, \ldots, N-1\}$ is the epistasis parameter.
  \item $\bA \in \{0,1\}^{N \times N}$ is the NK influence matrix with
        $A_{dd} = 1$ for all $d$ (each component depends on itself) and
        $\sum_{j} A_{dj} = K+1$ for all $d$.
  \item $f_d : \{0,1\}^{K+1} \to [0,1]$ is the component payoff function,
        with values drawn i.i.d.\ from $\mathrm{Uniform}(0,1)$.
\end{enumerate}
The fitness of a configuration $\bx \in \{0,1\}^N$ is
\begin{equation}
\label{eq:nk_fitness}
  W(\bx) = \frac{1}{N} \sum_{d=1}^{N} f_d\!\left(\bx_{\mathcal{N}(d)}\right),
\end{equation}
where $\bx_{\mathcal{N}(d)} = (x_j)_{j \in \mathcal{N}(d)}$ is the sub-vector of
$\bx$ restricted to the epistatic neighborhood $\mathcal{N}(d) = \{j : A_{dj}=1\}$.
\end{definition}

\begin{definition}[NK Adaptive Walk]
\label{def:nk_walk}
An \emph{adaptive walk} on an NK landscape is a sequence
$\bx^{(0)}, \bx^{(1)}, \ldots$ defined by:
\begin{enumerate}[label=(\roman*)]
  \item $\bx^{(0)}$ is drawn uniformly from $\{0,1\}^N$.
  \item At each step $t$, a component $d_t$ is selected (uniformly at random or
        sequentially), and the single-bit-flip
        $\bx^{(t)} \oplus \be_{d_t}$ is evaluated. If
        $W(\bx^{(t)} \oplus \be_{d_t}) > W(\bx^{(t)})$, then
        $\bx^{(t+1)} = \bx^{(t)} \oplus \be_{d_t}$; otherwise,
        $\bx^{(t+1)} = \bx^{(t)}$.
  \item The walk terminates at a local optimum where no single-bit flip
        improves fitness.
\end{enumerate}
A variant is the \emph{greedy adaptive walk}, where at each step the bit
flip yielding the largest fitness improvement is selected:
\begin{equation}
\label{eq:nk_greedy}
  d_t^* = \argmax_{d \in \{1,\ldots,N\}}
    \left[W\!\left(\bx^{(t)} \oplus \be_d\right) - W\!\left(\bx^{(t)}\right)\right].
\end{equation}
\end{definition}

% ============================================================
\subsection{The SaoMNK Model}
\label{sec:saomnk_model}
% ============================================================

\begin{definition}[SaoMNK Model]
\label{def:saomnk}
A \emph{Stochastic Actor-Oriented Model on NK Landscapes} (SaoMNK) is a tuple
$(M, N, \bE, \{f_d\}_{d=1}^N, \btheta, \beta)$ where:
\begin{enumerate}[label=(\roman*)]
  \item $M \in \N$ is the number of actors.
  \item $N \in \N$ is the number of binary components.
  \item $\bE \in \{0,1\}^{N \times N}$ is the influence matrix (with
        $E_{dd}=1$ for all $d$).
  \item $\{f_d\}_{d=1}^N$ is a set of component payoff functions, as in the NK
        model.
  \item $\btheta \in \R^p$ is the vector of SAOM effect parameters.
  \item $\beta > 0$ is the logit precision parameter.
\end{enumerate}
The state space is $\mathcal{S} = \{0,1\}^{M \times N}$, with generic element
$\bB$. The $i$-th row $\bb_i \in \{0,1\}^N$ represents actor $i$'s
configuration.
\end{definition}

\begin{definition}[SaoMNK Utility]
\label{def:saomnk_utility}
Actor $i$'s utility from configuration $\bb_i$, given the configurations of all
other actors $\bB_{-i}$ and influence matrix $\bE$, is
\begin{equation}
\label{eq:saomnk_utility}
  u_i(\bb_i; \bB_{-i}, \bE)
  = \underbrace{\sum_{d=1}^{N}
    f_d\!\left(\bb_i \odot \bE_{d\cdot}\right)}_{\text{NK component}}
  + \underbrace{\sum_{k=1}^{p} \theta_k \, s_k(\bb_i, \bB)}_{\text{SAOM network statistics}},
\end{equation}
where $\bE_{d\cdot}$ is the $d$-th row of $\bE$, the operation
$\bb_i \odot \bE_{d\cdot}$ denotes element-wise multiplication (masking
$\bb_i$ by the epistatic neighborhood of $d$), and $f_d$ is evaluated by
converting the masked binary sub-vector to an integer index via the power key
$\mathrm{PK}(N)$:
\begin{equation}
\label{eq:powerkey_index}
  f_d\!\left(\bb_i \odot \bE_{d\cdot}\right)
  \;=\; \mathrm{NK\_land}\!\left[\,
    \sum_{j=1}^{N} b_{ij}\, E_{dj}\, \mathrm{PK}(N)_j \;+\; 1,\; d\right].
\end{equation}
The SAOM statistics $s_k$ include network-structural effects defined on the
bipartite matrix $\bB$:
\begin{align}
  s_{\mathrm{density}}(\bb_i) &= \sum_{j=1}^N b_{ij}, \label{eq:density}\\
  s_{\mathrm{outAct}}(\bb_i) &= \left(\sum_{j=1}^N b_{ij}\right)^{\!2}, \label{eq:outact}\\
  s_{\mathrm{inPop}}(\bb_i, \bB) &= \sum_{j=1}^N b_{ij}\left(\sum_{m=1}^M b_{mj} + 1\right), \label{eq:inpop}\\
  s_{\mathrm{cycle4}}(\bb_i, \bB) &= \tfrac{1}{2}\,\bigl[\mathrm{diag}(\bB\bB^\top)^2 \cdot \bB\bB^\top\bigr]_{ii}. \label{eq:cycle4}
\end{align}
\end{definition}

\begin{definition}[SaoMNK Dynamics]
\label{def:saomnk_dynamics}
The SaoMNK model evolves as a continuous-time Markov chain of \emph{ministeps}.
At each ministep:
\begin{enumerate}[label=(\roman*)]
  \item An actor $i$ is selected according to a rate function $\lambda_i > 0$.
  \item Actor $i$ evaluates all possible single-tie changes
        $\bb_i \oplus \be_j$ for $j \in \{1,\ldots,N\}$, as well as maintaining
        the status quo $\bb_i$.
  \item The choice follows a multinomial logit:
        \begin{equation}
        \label{eq:logit_choice}
          P(\text{actor } i \text{ changes tie to } j)
          = \frac{\exp\!\bigl(\beta\, \Delta u_{ij}\bigr)}
                 {\displaystyle 1 + \sum_{k=1}^{N}
                  \exp\!\bigl(\beta\, \Delta u_{ik}\bigr)},
        \end{equation}
        where $\Delta u_{ij} = u_i(\bb_i \oplus \be_j; \bB_{-i}, \bE) -
        u_i(\bb_i; \bB_{-i}, \bE)$ is the change in utility from flipping
        component $j$.
\end{enumerate}
\end{definition}

% ============================================================
\section{Theorem 1: Reduction}
\label{sec:reduction}
% ============================================================

\begin{theorem}[Reduction Theorem]
\label{thm:reduction}
The standard NK model \emph{(}\Cref{def:nk,def:nk_walk}\emph{)} is a degenerate
special case of the SaoMNK model \emph{(}\Cref{def:saomnk,def:saomnk_utility,%
def:saomnk_dynamics}\emph{)}. Specifically, the NK landscape and its adaptive
walk are recovered exactly from SaoMNK under the following restriction:
\begin{enumerate}[label=(R\arabic*)]
  \item\label{R1} $M = 1$ (single actor);
  \item\label{R2} $\bE = \bA$ (SaoMNK influence matrix equals the NK influence
        matrix, with $K+1$ ones per row);
  \item\label{R3} $\btheta = \mathbf{0}$ (all SAOM network-statistic parameters
        set to zero);
  \item\label{R4} $\beta \to \infty$ (logit precision diverges, recovering
        deterministic greedy choice);
  \item\label{R5} $\lambda_i$ constant across components (uniform selection
        probability).
\end{enumerate}
\end{theorem}

\begin{proof}
We establish the result in two parts: payoff equivalence
(\Cref{lem:payoff_equiv}) and dynamic equivalence
(\Cref{lem:dynamic_equiv}).

\begin{lemma}[Payoff Equivalence]
\label{lem:payoff_equiv}
Under restrictions \ref{R1}--\ref{R3}, the SaoMNK utility function reduces to
the NK fitness function (up to a constant scaling factor of $N$).
\end{lemma}

\begin{proof}[Proof of \Cref{lem:payoff_equiv}]
Under \ref{R1}, there is a single actor ($M=1$), so the bipartite matrix
$\bB$ reduces to a single row vector $\bb_1 \equiv \bx \in \{0,1\}^N$. The
full state space collapses from $\{0,1\}^{M \times N}$ to $\{0,1\}^N$.

Under \ref{R3}, all SAOM network-statistic parameters are zero
($\btheta = \mathbf{0}$), so the network component in
\eqref{eq:saomnk_utility} vanishes:
\begin{equation*}
  \sum_{k=1}^{p} \theta_k \, s_k(\bb_1, \bB) = 0.
\end{equation*}
Moreover, with $M = 1$ and $\btheta = \mathbf{0}$, the statistics
$s_{\mathrm{inPop}}$ and $s_{\mathrm{cycle4}}$---which depend on
$\bB_{-i}$---become irrelevant.

The SaoMNK utility \eqref{eq:saomnk_utility} therefore simplifies to
\begin{equation}
\label{eq:reduced_utility}
  u(\bx; \bE) = \sum_{d=1}^{N} f_d\!\left(\bx \odot \bE_{d\cdot}\right).
\end{equation}

Under \ref{R2}, $\bE = \bA$, so $\bE_{d\cdot} = \bA_{d\cdot}$ and the
mask $\bx \odot \bE_{d\cdot}$ selects exactly the entries of $\bx$ in the
epistatic neighborhood $\mathcal{N}(d) = \{j : A_{dj} = 1\}$.

Under \ref{R2}, every term of \eqref{eq:reduced_utility} is the NK component
payoff evaluated on the same sub-vector, so
\begin{equation}
\label{eq:payoff_equiv}
  u(\bx; \bA) = \sum_{d=1}^{N} f_d(\bx_{\mathcal{N}(d)}) = N \cdot W(\bx)
  \quad\text{for every } \bx \in \{0,1\}^N .
\end{equation}
Every component contributes its payoff whether $x_d = 0$ or $x_d = 1$, exactly as
in \Cref{def:nk}; this is why the NK term of \eqref{eq:saomnk_utility} carries
no factor $b_{id}$. (An earlier version of this proof multiplied each payoff by
$b_{id}$ and verified the identity only at $\bx = \mathbf{1}$; with that factor
the identity fails whenever some $x_d = 0$, and the greedy walks can move in
opposite directions from the same state. The counterexample and the repaired
identity are machine-checked in Lean 4; see the note at the head of this file.)
The factor $\frac{1}{N}$ is a normalization constant that does not affect
ordinal rankings of configurations.

This completes the payoff equivalence: the SaoMNK utility under
\ref{R1}--\ref{R3} is an affine transformation of the NK fitness function, and
all ordinal comparisons---which drive adaptive search---are preserved.
\end{proof}

\begin{lemma}[Dynamic Equivalence]
\label{lem:dynamic_equiv}
Under restrictions \ref{R1}--\ref{R5}, the SaoMNK ministep dynamics produce the
same sequence of states as the NK greedy adaptive walk.
\end{lemma}

\begin{proof}[Proof of \Cref{lem:dynamic_equiv}]
Under \ref{R1}, there is only one actor, so the rate function is trivial: the
single actor is always selected. Under \ref{R5}, each component is equally
likely to be the focal tie for evaluation.

At each ministep, the actor evaluates all $N$ possible single-tie changes.
The change in utility from flipping component $j$ is
\begin{equation*}
  \Delta u_j = u(\bx \oplus \be_j; \bA) - u(\bx; \bA).
\end{equation*}
By \Cref{lem:payoff_equiv}, this is proportional to the change in NK fitness:
$\Delta u_j = N \cdot [W(\bx \oplus \be_j) - W(\bx)]$.

The SaoMNK logit choice probability \eqref{eq:logit_choice} is
\begin{equation*}
  P(\text{flip } j) = \frac{\exp(\beta \, \Delta u_j)}{1 + \sum_{k=1}^{N} \exp(\beta \, \Delta u_k)}.
\end{equation*}

Now take the limit $\beta \to \infty$ (\ref{R4}). Let
$j^* = \argmax_{j} \Delta u_j$ and assume the maximum is unique (which holds
with probability 1 under i.i.d.\ uniform payoffs). Then:
\begin{equation*}
  \lim_{\beta \to \infty} P(\text{flip } j)
  = \begin{cases}
    1 & \text{if } j = j^* \text{ and } \Delta u_{j^*} > 0, \\
    0 & \text{otherwise}.
  \end{cases}
\end{equation*}

If $\Delta u_{j^*} > 0$, the actor deterministically selects the utility-maximizing
flip. If $\Delta u_j \leq 0$ for all $j$, then the status quo probability
$1/(1 + \sum_k \exp(\beta \Delta u_k))$ converges to 1, and the process remains
at $\bx$---a local optimum.

This is exactly the NK greedy adaptive walk (\Cref{def:nk_walk}, equation
\eqref{eq:nk_greedy}): at each step, the bit flip yielding the largest fitness
improvement is selected, and the walk terminates at a local optimum.

Under \ref{R5} with uniform component selection (the ``random'' variant of the
adaptive walk), the reduction yields the random-order adaptive walk rather than
the greedy walk, but both are standard NK dynamics.
\end{proof}

Combining \Cref{lem:payoff_equiv} and \Cref{lem:dynamic_equiv}, we conclude
that the NK model---both its fitness landscape and its adaptive walk
dynamics---is recovered exactly from SaoMNK under restrictions
\ref{R1}--\ref{R5}. Therefore, the standard NK model is a degenerate special
case of SaoMNK: $\mathrm{NK} \subset \mathrm{SaoMNK}$. \qedhere
\end{proof}

\begin{remark}[Computational Implementation]
\label{rem:implementation}
The payoff equivalence can be verified directly in the SaoMNK codebase.
The function \texttt{calc\_fit()} computes fitness via the expression
\begin{center}
\texttt{NK\_land[sum(Current\_position * inter\_m[d, ] * Power\_key) + 1, d]}
\end{center}
which implements \eqref{eq:powerkey_index}. Setting \texttt{inter\_m} $= \bA$
and \texttt{Current\_position} $= \bx$, this is exactly the NK payoff lookup
$f_d(\bx_{\mathcal{N}(d)})$.
\end{remark}

% ============================================================
\section{Theorem 2: Strict Generalization}
\label{sec:generalization}
% ============================================================

\begin{theorem}[Generalization Theorem]
\label{thm:generalization}
SaoMNK strictly generalizes the NK model along three independent dimensions:
\begin{enumerate}[label=(\alph*)]
  \item \emph{Multi-actor endogenous interaction}: $M > 1$ actors jointly
        navigate a shared landscape.
  \item \emph{Endogenous landscape co-evolution}: each actor's effective
        landscape depends on the current configurations of other actors
        through SAOM network statistics.
  \item \emph{Bounded rationality with tunable search precision}: the logit
        parameter $\beta \in (0,\infty)$ interpolates continuously between
        random and deterministic search.
\end{enumerate}
Each dimension independently produces behavior that cannot be represented in
any standard NK model, regardless of the choice of $N$ or $K$.
\end{theorem}

\begin{proof}
We establish each dimension in turn.

\medskip
\noindent\textbf{Dimension (a): Multi-actor state space.}

The NK model operates on a single binary string $\bx \in \{0,1\}^N$, yielding
a state space of cardinality $|\mathcal{S}_{\mathrm{NK}}| = 2^N$. The SaoMNK
model, with $M$ actors, operates on $\bB \in \{0,1\}^{M \times N}$, yielding
$|\mathcal{S}_{\mathrm{SaoMNK}}| = 2^{MN}$.

For $M \geq 2$, the state space grows exponentially:
$2^{MN} = (2^N)^M$. But the significance extends beyond cardinality. In SaoMNK,
the $M$ actors share the same $N$ components and payoff table $\{f_d\}$, but
each actor independently controls a row of $\bB$. This creates a \emph{game}
among actors over the landscape, where interactions arise endogenously through
the shared component space.

In contrast, the standard NK model is inherently a single-agent optimization
problem. Even when multiple agents are introduced \emph{ad hoc} (as in
\citealt{siggelkow2005speed}), they operate on \emph{disjoint} subsets of
components with exogenously specified cross-agent coupling. SaoMNK provides
a principled micro-foundation for multi-agent search via the SAOM
framework \citep{snijders2010introduction}.

\medskip
\noindent\textbf{Dimension (b): Endogenous landscape co-evolution.}

This is the most consequential departure from the NK model. In the standard
NK model (and its multi-agent extensions), the fitness landscape $W(\bx)$
is fixed and exogenous: it depends only on the payoff table $\{f_d\}$ and
the NK influence matrix $\bA$. Other agents' configurations do not enter the
payoff function.

In SaoMNK, the SAOM network statistics in \eqref{eq:saomnk_utility}
create dependencies between actors. Consider the \emph{inPop} statistic
\eqref{eq:inpop}:
\begin{equation*}
  s_{\mathrm{inPop}}(\bb_i, \bB) = \sum_{j=1}^{N} b_{ij}
    \left(\sum_{m=1}^{M} b_{mj} + 1\right).
\end{equation*}
Actor $i$'s utility from connecting to component $j$ depends on
$\sum_{m} b_{mj}$, the number of \emph{other} actors also connected to $j$.
If $\theta_{\mathrm{inPop}} > 0$, actors prefer popular components
(agglomeration); if $\theta_{\mathrm{inPop}} < 0$, actors avoid
crowded components (differentiation).

\begin{lemma}[Landscape Endogeneity]
\label{lem:landscape_endogeneity}
When $M \geq 2$ and $\theta_{\mathrm{inPop}} \neq 0$, actor $i$'s effective
fitness landscape $\tilde{W}_i(\bb_i) \triangleq u_i(\bb_i; \bB_{-i}, \bE)$
depends on $\bB_{-i}$. Hence, any change in actor $j$'s configuration
($j \neq i$) generically alters actor $i$'s landscape: for all $j \neq i$,
there exist configurations $\bb_j, \bb_j'$ such that
\begin{equation*}
  \tilde{W}_i(\bb_i \mid \bb_j) \neq \tilde{W}_i(\bb_i \mid \bb_j')
  \quad\text{for some } \bb_i.
\end{equation*}
\end{lemma}

\begin{proof}[Proof of \Cref{lem:landscape_endogeneity}]
Fix actor $i$ with configuration $\bb_i$ where $b_{id} = 1$ for some
component $d$. Consider actor $j \neq i$ with two configurations
$\bb_j$ and $\bb_j'$ that differ only in component $d$: $b_{jd} = 0$ in
$\bb_j$ and $b_{jd}' = 1$ in $\bb_j'$. Then:
\begin{equation*}
  s_{\mathrm{inPop}}(\bb_i, \bB') - s_{\mathrm{inPop}}(\bb_i, \bB)
  = b_{id} \cdot 1 = 1 \neq 0,
\end{equation*}
so $u_i(\bb_i; \bB_{-i}', \bE) - u_i(\bb_i; \bB_{-i}, \bE) =
\theta_{\mathrm{inPop}} \neq 0$.
\end{proof}

This endogenous landscape co-evolution is fundamentally different from
``coupled landscapes'' \citep{siggelkow2005speed, rivkin2007balancing},
where inter-agent coupling is specified exogenously via a fixed matrix. In
SaoMNK, the coupling arises endogenously through the \emph{evolving}
network structure $\bB$.

The \emph{cycle4} statistic \eqref{eq:cycle4} creates even richer
dependencies: it captures four-cycle structures in the bipartite graph,
meaning that actor $i$'s utility depends on higher-order patterns of
co-selection among multiple actors and components.

\medskip
\noindent\textbf{Dimension (c): Bounded rationality.}

In the NK model, search is deterministic: the adaptive walk always selects the
best-improving single-bit flip (greedy) or any improving flip (random). There
is no mechanism for sub-optimal choices, exploration, or error.

SaoMNK uses the logit choice model \eqref{eq:logit_choice}, parameterized by
$\beta > 0$. This yields a continuous spectrum of search behaviors:
\begin{itemize}
  \item $\beta \to 0$: Each alternative is chosen with equal probability
        (random search). The actor ignores payoff differences entirely.
  \item $\beta \in (0, \infty)$: Noisy best-response. The actor tends toward
        utility-maximizing choices but makes errors with probability
        decreasing in $\beta$ and in the utility gap. This corresponds to
        bounded rationality in the sense of \citet{blume1993statistical}
        and the conditional logit of \citet{mcfadden1974conditional}.
  \item $\beta \to \infty$: Deterministic greedy choice (the NK limit, as
        shown in \Cref{lem:dynamic_equiv}).
\end{itemize}

\begin{lemma}[Search Precision Spectrum]
\label{lem:search_precision}
For any $\beta \in (0, \infty)$ and any configuration $\bx$ that is not a
global optimum, the SaoMNK dynamics with $M=1$ and $\btheta = \mathbf{0}$
assign positive probability to \emph{every} single-bit flip, including
fitness-decreasing ones. In contrast, the NK adaptive walk assigns zero
probability to fitness-decreasing flips.
\end{lemma}

\begin{proof}[Proof of \Cref{lem:search_precision}]
For any $\beta \in (0, \infty)$ and any $j \in \{1, \ldots, N\}$:
\begin{equation*}
  P(\text{flip } j) = \frac{\exp(\beta \, \Delta u_j)}{1 + \sum_{k=1}^{N} \exp(\beta \, \Delta u_k)} > 0,
\end{equation*}
since $\exp(\beta \, \Delta u_j) > 0$ for all finite $\beta$ and all
$\Delta u_j \in \R$, and the denominator is finite. In particular, even if
$\Delta u_j < 0$ (fitness-decreasing flip), the probability is strictly
positive. The NK adaptive walk, by contrast, assigns probability 0 to any
flip with $\Delta u_j \leq 0$.
\end{proof}

This has important substantive implications. With finite $\beta$, actors can
escape local optima---a behavior impossible in standard NK. The parameter
$\beta$ therefore controls the exploration--exploitation tradeoff, a central
concern in organizational learning theory.

\medskip
\noindent\textbf{Strict generalization.}

Each dimension produces behavior outside the NK model family independently:

\begin{example}[Dimension (a) alone]
With $M = 2$, $\btheta = \mathbf{0}$, and $\beta \to \infty$, SaoMNK
reduces to two independent NK adaptive walks on the same landscape. While
trivial, this already lies outside the single-agent NK framework.
\end{example}

\begin{example}[Dimension (b) alone]
With $M = 2$, $\theta_{\mathrm{inPop}} = 1$, and $\beta \to \infty$, actor
1's landscape changes whenever actor 2 changes a tie. Specifically, if actor
2 connects to component $d$, then actor 1's utility from connecting to $d$
increases by $\theta_{\mathrm{inPop}} = 1$. This dynamic landscape
reshaping cannot be captured by \emph{any} single-actor NK model with any
$N$ or $K$, because NK fitness $W(\bx)$ depends only on $\bx$ and the fixed
payoff table---it has no mechanism for another agent's state to enter the
payoff function.
\end{example}

\begin{example}[Dimension (c) alone]
With $M = 1$, $\btheta = \mathbf{0}$, and $\beta = 1$, SaoMNK on the same
landscape as an NK model produces a stochastic search trajectory that
assigns positive probability to fitness-decreasing moves
(\Cref{lem:search_precision}). No parameterization of the NK adaptive walk
can reproduce this: NK search is deterministic.
\end{example}

Since each dimension independently generates behavior impossible in any NK
model, and \Cref{thm:reduction} shows that NK is a special case of SaoMNK,
the generalization is strict: $\mathrm{NK} \subsetneq \mathrm{SaoMNK}$.
\end{proof}

% ============================================================
\section{Theorem 3: Approximation}
\label{sec:approximation}
% ============================================================

\begin{theorem}[Approximation Theorem]
\label{thm:approximation}
For any NK landscape $(N, K, \bA, \{f_d\}_{d=1}^N)$, there exists a
parameterization of the SaoMNK model that reproduces the NK payoff structure to
arbitrary precision.
\end{theorem}

We establish this result via three independent proof strategies, each
illuminating a different aspect of the approximation.

% ---------- 3a ----------
\begin{proposition}[Constructive Existence via Dummy Covariates]
\label{prop:constructive}
For any NK landscape $(N, K, \bA, \{f_d\}_{d=1}^N)$, there exists a set of
$N \cdot 2^{K+1}$ dyadic covariates and corresponding SaoMNK parameters that
\emph{exactly} reproduce the NK fitness function.
\end{proposition}

\begin{proof}
Fix a component $d$ and its epistatic neighborhood $\mathcal{N}(d)$ with
$|\mathcal{N}(d)| = K + 1$. There are $2^{K+1}$ possible configurations of
the bits in $\mathcal{N}(d)$. Enumerate these configurations as
$\mathbf{c}_1^{(d)}, \mathbf{c}_2^{(d)}, \ldots, \mathbf{c}_{2^{K+1}}^{(d)}$.

For each component $d$ and each configuration $\mathbf{c}_\ell^{(d)}$, define
a dummy-coded dyadic covariate:
\begin{equation}
\label{eq:dummy_covariate}
  w_{d\ell}(\bx) =
  \begin{cases}
    1 & \text{if } \bx_{\mathcal{N}(d)} = \mathbf{c}_\ell^{(d)}, \\
    0 & \text{otherwise}.
  \end{cases}
\end{equation}
This is a \emph{dyadic covariate} in the SAOM sense: it depends on the
actor's tie to component $d$ and on the configuration of related
components.

Set the corresponding SaoMNK parameter equal to the NK payoff value:
$\alpha_{d\ell} = f_d(\mathbf{c}_\ell^{(d)})$.

Then the NK component payoff is recovered exactly:
\begin{equation*}
  f_d(\bx_{\mathcal{N}(d)}) = \sum_{\ell=1}^{2^{K+1}} \alpha_{d\ell} \, w_{d\ell}(\bx),
\end{equation*}
since exactly one dummy is active ($w_{d\ell} = 1$) for each configuration,
and it selects the correct payoff value.

Summing over all $N$ components:
\begin{equation*}
  \sum_{d=1}^{N} f_d(\bx_{\mathcal{N}(d)})
  = \sum_{d=1}^{N} \sum_{\ell=1}^{2^{K+1}} \alpha_{d\ell} \, w_{d\ell}(\bx)
  = N \cdot W(\bx).
\end{equation*}

This construction uses $N \cdot 2^{K+1}$ covariates and parameters, which
grows exponentially in $K$. The construction is therefore impractical for
large $K$ but serves as an \emph{existence proof}: the SaoMNK parameter
space is rich enough to exactly encode any NK landscape.
\end{proof}

\begin{remark}
The exponential scaling in \Cref{prop:constructive} mirrors the fundamental
dimensionality of the NK model itself: the payoff table $\{f_d\}$ contains
$N \cdot 2^{K+1}$ independent entries. No representation can be more compact
without imposing structure on the payoffs.
\end{remark}

% ---------- 3b ----------
\begin{proposition}[Actor Heterogeneity Substitution]
\label{prop:heterogeneity}
With $M$ heterogeneous actors and $Q$ continuous dyadic covariates interacted
with the epistasis structure, the SaoMNK model has $Q \cdot M \cdot N$
effective parameters. A sufficient condition for matching the NK payoff
dimensionality is $Q \cdot M \geq 2^{K+1}$.
\end{proposition}

\begin{proof}
In the SaoMNK model, let $\{z_q^{(i,d)}\}_{q=1}^Q$ denote $Q$ continuous
dyadic covariates for actor $i$ and component $d$. These covariates enter the
utility function through SAOM effects of the form \texttt{X}, \texttt{egoX},
\texttt{altX}, and \texttt{XWX}:
\begin{equation}
\label{eq:covariate_utility}
  u_i^{\mathrm{cov}}(\bb_i) = \sum_{d=1}^{N} b_{id} \sum_{q=1}^{Q}
    \gamma_q \, z_q^{(i,d)},
\end{equation}
where $\gamma_q$ are the covariate effect parameters.

Since the covariates can vary across both actors $i \in \{1, \ldots, M\}$
and components $d \in \{1, \ldots, N\}$, the total number of free covariate
values is $Q \cdot M \cdot N$.

The NK payoff table has $N \cdot 2^{K+1}$ independent entries (each of the
$N$ components has a payoff function over $2^{K+1}$ configurations). For the
SaoMNK model to have at least as many degrees of freedom, we require:
\begin{equation}
\label{eq:heterogeneity_condition}
  Q \cdot M \geq 2^{K+1}.
\end{equation}

This condition is practically achievable. Consider the following examples:
\begin{itemize}
  \item $M = 10$, $K = 3$: Need $Q \cdot 10 \geq 16$, so $Q = 2$ suffices.
  \item $M = 10$, $K = 5$: Need $Q \cdot 10 \geq 64$, so $Q = 7$ suffices.
  \item $M = 20$, $K = 7$: Need $Q \cdot 20 \geq 256$, so $Q = 13$ suffices.
  \item $M = 50$, $K = 10$: Need $Q \cdot 50 \geq 2048$, so $Q = 41$ suffices.
\end{itemize}

The key insight is that \emph{actor heterogeneity provides a natural,
empirically grounded source of landscape complexity} that substitutes for
the purely random payoffs of the NK model. Rather than drawing $2^{K+1}$
random numbers per component, the SaoMNK model generates effective landscape
ruggedness through the interaction of actor-specific attributes with the
epistatic structure.

To make the dimension count precise: fix component $d$. In the NK model,
$f_d$ maps $2^{K+1}$ configurations to $[0,1]$, requiring $2^{K+1}$ parameters.
In SaoMNK, for component $d$, there are $Q$ covariates $\times$ $M$ actors
$= Q \cdot M$ actor-specific payoff modifiers. When $Q \cdot M \geq 2^{K+1}$,
the linear span of these modifiers has sufficient dimension to approximate
any function on $\{0,1\}^{K+1}$ (by the argument in \Cref{prop:constructive},
replacing dummy covariates with continuous ones whose span covers the same
space).

This result suggests that models with moderate numbers of actors and covariates
can capture landscapes of substantial complexity---and can do so with
\emph{empirically estimable} parameters rather than arbitrary random draws.
\end{proof}

% ---------- 3c ----------
\begin{proposition}[Universal Approximation via Mixed Logit]
\label{prop:mixed_logit}
The SaoMNK model with random effects (mixed logit specification) can
approximate any NK payoff function to arbitrary precision.
\end{proposition}

\begin{proof}
This result follows from the universal approximation theorem for mixed logit
models established by \citet{mcfadden2000mixed}.

\medskip
\noindent\emph{Step 1: The SAOM is a conditional logit.}

The SaoMNK choice probability \eqref{eq:logit_choice} has the form
\begin{equation*}
  P(j \mid \bb_i, \bB_{-i}) = \frac{\exp(V_{ij})}{\sum_k \exp(V_{ik})},
\end{equation*}
where $V_{ij} = \beta \, \Delta u_{ij}$ is a deterministic function of
observable attributes of the alternatives (components) and the chooser
(actor). This is precisely the conditional logit model of
\citet{mcfadden1974conditional}.

\medskip
\noindent\emph{Step 2: The mixed logit extension.}

A \emph{mixed logit} (or random-parameters logit) generalizes the conditional
logit by allowing the parameters $\btheta$ to vary across decision-makers
according to a mixing distribution $G$:
\begin{equation}
\label{eq:mixed_logit}
  P(j \mid \bb_i, \bB_{-i}) = \int
    \frac{\exp(V_{ij}(\btheta))}{\sum_k \exp(V_{ik}(\btheta))}
    \, dG(\btheta).
\end{equation}

In the SaoMNK framework, this corresponds to allowing the effect parameters
$\btheta$ (and/or $\beta$) to be drawn from a distribution across actors
or across time periods, rather than being fixed.

\medskip
\noindent\emph{Step 3: The McFadden--Train theorem.}

\citet{mcfadden2000mixed} prove the following (Theorem~1 in their paper):

\begin{quote}
\emph{Any Random Utility Model (RUM) can be approximated to any degree of
accuracy by a mixed logit model with an appropriate choice of mixing
distribution and variables.}
\end{quote}

Formally, let $\mathcal{P}^*$ denote any choice probability system consistent
with a random utility model. For any $\varepsilon > 0$, there exists a
mixed logit specification---i.e., a set of alternative-specific variables
and a mixing distribution $G$ over parameters---such that
\begin{equation*}
  \sup_j \left|P(j) - P^*(j)\right| < \varepsilon.
\end{equation*}

\medskip
\noindent\emph{Step 4: Application to NK payoffs.}

The NK fitness function $W(\bx)$ defines a preference ordering over
configurations $\bx \in \{0,1\}^N$. This preference ordering can be
rationalized by a random utility model: set $U(\bx) = W(\bx) + \eta$
where $\eta$ is a noise term.

By the McFadden--Train theorem, the choice probabilities implied by this
RUM can be approximated to arbitrary precision by a mixed logit. Since the
SAOM is a conditional logit, its mixed version (with random effects on
$\btheta$) is a mixed logit. Therefore, the mixed SaoMNK model can
approximate the choice behavior induced by any NK fitness function to
arbitrary precision.

Moreover, since the approximation holds for choice probabilities, and
the choice probabilities under the logit model are monotone transformations
of utility differences, the underlying \emph{payoff structure} is also
approximated: for any two configurations $\bx, \bx'$ with
$W(\bx) > W(\bx')$, the mixed SaoMNK model can be parameterized so that
$u(\bx) > u(\bx')$ with probability arbitrarily close to 1.
\end{proof}

\begin{remark}
\Cref{prop:mixed_logit} is the strongest theoretical result: it provides
a \emph{universal} approximation guarantee without requiring exponentially
many covariates (\Cref{prop:constructive}) or large numbers of actors
(\Cref{prop:heterogeneity}). The cost is that the mixing distribution
$G$ may be complex, but it can in principle be estimated from data using
Bayesian methods or simulated maximum likelihood
\citep{mcfadden2000mixed}.
\end{remark}

\begin{corollary}
\label{cor:approximation_hierarchy}
The three approximation strategies form a hierarchy of increasing
generality:
\begin{enumerate}[label=(\roman*)]
  \item \Cref{prop:constructive} provides exact reconstruction but
        requires $O(N \cdot 2^K)$ parameters.
  \item \Cref{prop:heterogeneity} provides exact reconstruction using
        $O(Q \cdot M \cdot N)$ empirically meaningful parameters, where
        $Q \cdot M \geq 2^{K+1}$.
  \item \Cref{prop:mixed_logit} provides $\varepsilon$-approximation for
        any $\varepsilon > 0$ with a flexible nonparametric mixing
        distribution, and requires no explosion in the number of observed
        covariates.
\end{enumerate}
\end{corollary}

% ============================================================
\section{Discussion}
\label{sec:discussion}
% ============================================================

The three theorems established above---Reduction, Generalization, and
Approximation---collectively position the SaoMNK model as a principled
generalization of the NK landscape framework that has shaped complexity-theoretic
research in strategic management for three decades
\citep{kauffman1993origins, levinthal1997adaptation, rivkin2007balancing,
siggelkow2005speed}. We discuss the implications for theory, methodology,
and future research.

\subsection{Theoretical Implications}

\paragraph{From landscapes to games.}
The Reduction Theorem (\Cref{thm:reduction}) shows that the NK model is a
single-actor, deterministic-search, fixed-landscape special case of SaoMNK.
This reframing is consequential: it reveals that the NK model implicitly
assumes away three features---multi-actor interaction, landscape endogeneity,
and bounded rationality---that are central to organizational and strategic
phenomena.

The Generalization Theorem (\Cref{thm:generalization}) demonstrates that
relaxing any one of these assumptions produces behavior outside the NK family.
Dimension (b)---endogenous landscape co-evolution---is particularly important.
Prior multi-agent extensions of NK \citep{siggelkow2005speed, rivkin2007balancing}
couple agents' landscapes exogenously, specifying in advance which of agent
$j$'s decisions affect agent $i$'s fitness. In SaoMNK, this coupling emerges
endogenously from the evolving network structure: when actor $j$ connects to
(or disconnects from) a component, it changes the popularity, crowding, and
structural statistics that enter actor $i$'s utility function. The landscape
is not merely ``coupled''---it \emph{co-evolves} with the actors' search
trajectories.

\paragraph{Bounded rationality as a first-class citizen.}
The logit choice model (\Cref{lem:search_precision}) provides a natural,
axiomatically grounded model of bounded rationality that the NK framework
lacks. The precision parameter $\beta$ has a clear behavioral interpretation:
it governs the probability that an actor deviates from the payoff-maximizing
choice. When $\beta$ is small, actors explore broadly (high error rates);
when $\beta$ is large, they exploit locally (low error rates). This connects
the complexity-theoretic tradition to the behavioral theory of the firm and
to evolutionary game theory \citep{blume1993statistical}, where logit dynamics
are a standard model of noisy best-response.

The ability to escape local optima under finite $\beta$ is substantively
important. A large literature has documented the ``competency trap'' and
related phenomena in which organizations become locked into suboptimal
configurations. The NK model can represent such lock-in (local optima on
rugged landscapes) but has no mechanism for escape. SaoMNK, with finite
$\beta$, naturally generates both lock-in \emph{and} escape, with the
relative prevalence governed by a single estimable parameter.

\subsection{Methodological Implications}

\paragraph{From simulation to estimation.}
Perhaps the most significant practical advantage of SaoMNK over NK is
\emph{empirical estimability}. The NK model's payoff tables are drawn from
arbitrary random distributions and have no empirical counterpart: there is
no procedure for estimating $f_d$ from data. The model is therefore
restricted to computational experiments (simulation).

SaoMNK, by contrast, inherits the full statistical infrastructure of the
Stochastic Actor-Oriented Model \citep{snijders2010introduction}. The
parameters $\btheta$ and $\beta$ can be estimated from longitudinal
network data using the method of moments or maximum likelihood, as
implemented in the RSiena package. This opens the door to
\emph{empirical complexity theory}: researchers can estimate the
parameters governing search, interaction, and landscape structure from
real organizational data, rather than assuming them.

The Approximation Theorem (\Cref{thm:approximation}) provides the
theoretical foundation for this program. \Cref{prop:heterogeneity} shows
that actor heterogeneity---a feature naturally present in empirical
settings---substitutes for the arbitrary randomness of NK payoffs.
\Cref{prop:mixed_logit} guarantees that, with sufficient flexibility in
the mixing distribution, any NK-like payoff structure can be captured.

\paragraph{Parameter economy and interpretability.}
\Cref{prop:heterogeneity} highlights an important practical advantage:
with $M = 10$ actors and $Q = 2$ continuous covariates, SaoMNK can match
the complexity of an $N$-component, $K = 3$ NK landscape. The SaoMNK
parameters ($\theta_{\mathrm{density}}$, $\theta_{\mathrm{inPop}}$,
$\theta_{\mathrm{cycle4}}$, and covariate effects) have clear substantive
interpretations---they describe preferences for scope, popularity,
clustering, and covariate-driven matching. The NK payoff table, by
contrast, has $N \cdot 2^{K+1}$ uninterpretable random numbers.

\subsection{Connections to the Broader Literature}

The relationship between SaoMNK and NK connects several previously
disparate literatures:
\begin{itemize}
  \item \emph{Complexity theory in management}
        \citep{kauffman1993origins, levinthal1997adaptation}: SaoMNK
        preserves the core insight---that epistatic interactions among
        components create rugged fitness landscapes with multiple local
        optima---while adding micro-founded behavioral and social structure.

  \item \emph{Network dynamics}
        \citep{snijders2010introduction}: SaoMNK applies the SAOM
        framework, originally developed for social networks, to the
        domain of organizational search. The bipartite network of actors
        and components provides a natural representation of resource
        allocation, technology adoption, and market entry decisions.

  \item \emph{Discrete choice theory}
        \citep{mcfadden1974conditional, mcfadden2000mixed}: The logit
        choice model in SaoMNK connects organizational search to the
        econometric tradition of random utility maximization, enabling
        welfare analysis and counterfactual prediction.

  \item \emph{Evolutionary game theory}
        \citep{blume1993statistical}: The logit dynamics governing SaoMNK
        search are a form of stochastic best-response dynamics, linking
        the model to equilibrium selection in games.
\end{itemize}

\subsection{Limitations and Boundary Conditions}

Three limitations merit discussion. First, the Approximation Theorem
(\Cref{thm:approximation}) is an existence result: it guarantees that a
parameterization exists but does not provide an efficient algorithm for
finding it. In practice, estimation via RSiena's method of moments is
computationally feasible for moderate $M$ and $N$ but may become
challenging for very large systems.

Second, the Generalization Theorem (\Cref{thm:generalization}) shows that
SaoMNK is \emph{strictly} more expressive than NK, but greater
expressiveness carries the risk of overfitting. Empirical applications
should employ standard model selection techniques (BIC, cross-validation)
to identify the most parsimonious specification.

Third, the relationship established here is between SaoMNK and the
\emph{standard} NK model. Several extensions of NK---including NKC
\citep{kauffman1993origins}, NK with changing landscapes, and NK with
imperfect information---address some of the same limitations that SaoMNK
addresses. The formal relationship between SaoMNK and these extended NK
models remains an open question for future work.

\subsection{Conclusion}

The standard NK model has been a productive tool for generating
theoretical insights about the structure of complex adaptive systems. By
proving that it is a degenerate special case of SaoMNK---recoverable by
setting $M = 1$, $\btheta = \mathbf{0}$, and $\beta \to \infty$---we
provide a precise characterization of the assumptions implicit in NK-based
research. The Generalization Theorem identifies exactly how SaoMNK extends
beyond NK, and the Approximation Theorem establishes that this extension
does not sacrifice any of NK's representational power. Together, these
results provide a formal foundation for moving from simulation-only
complexity models to empirically estimable models of organizational search
on rugged landscapes.

% ============================================================
% REFERENCES
% ============================================================

\bibliographystyle{apalike}

\begin{thebibliography}{99}

\bibitem[Altenberg, 1997]{altenberg1997nk}
Altenberg, L. (1997).
\newblock {NK} fitness landscapes.
\newblock In B\"ack, T., Fogel, D., and Michalewicz, Z., editors,
  \emph{Handbook of Evolutionary Computation}, section B2.7.2.
  Oxford University Press.

\bibitem[Blume, 1993]{blume1993statistical}
Blume, L.~E. (1993).
\newblock The statistical mechanics of strategic interaction.
\newblock \emph{Games and Economic Behavior}, 5(3):387--424.

\bibitem[Kauffman, 1993]{kauffman1993origins}
Kauffman, S.~A. (1993).
\newblock \emph{The Origins of Order: Self-Organization and Selection in
  Evolution}.
\newblock Oxford University Press, New York.

\bibitem[Levinthal, 1997]{levinthal1997adaptation}
Levinthal, D.~A. (1997).
\newblock Adaptation on rugged landscapes.
\newblock \emph{Management Science}, 43(7):934--950.

\bibitem[McFadden, 1974]{mcfadden1974conditional}
McFadden, D. (1974).
\newblock Conditional logit analysis of qualitative choice behavior.
\newblock In Zarembka, P., editor, \emph{Frontiers in Econometrics},
  pages 105--142. Academic Press, New York.

\bibitem[McFadden and Train, 2000]{mcfadden2000mixed}
McFadden, D. and Train, K. (2000).
\newblock Mixed {MNL} models for discrete response.
\newblock \emph{Journal of Applied Econometrics}, 15(5):447--470.

\bibitem[Rivkin and Siggelkow, 2007]{rivkin2007balancing}
Rivkin, J.~W. and Siggelkow, N. (2007).
\newblock Patterned interactions in complex systems: Implications for
  exploration.
\newblock \emph{Management Science}, 53(7):1068--1085.

\bibitem[Siggelkow and Rivkin, 2005]{siggelkow2005speed}
Siggelkow, N. and Rivkin, J.~W. (2005).
\newblock Speed and search: Designing organizations for turbulence and
  complexity.
\newblock \emph{Organization Science}, 16(2):101--122.

\bibitem[Snijders et~al., 2010]{snijders2010introduction}
Snijders, T.~A.~B., van~de Bunt, G.~G., and Steglich, C.~E.~G. (2010).
\newblock Introduction to stochastic actor-based models for network dynamics.
\newblock \emph{Social Networks}, 32(1):44--60.

\end{thebibliography}

\end{document}
| \textit{(after
removing the duplicate preamble).}

\bigskip\noindent
\textbf{Summary of results:}

\begin{enumerate}
  \item \textbf{Theorem 1 (Reduction):} Standard NK is recovered from
    SaoMNK by setting $M=1$, $\boldsymbol{\theta}=\mathbf{0}$,
    $\beta\to\infty$, $\mathbf{E}=\mathbf{A}$, and constant rate functions.
    Proved via Payoff Equivalence (Lemma 1) and Dynamic Equivalence (Lemma 2).

  \item \textbf{Theorem 2 (Generalization):} SaoMNK strictly extends NK
    along three independent dimensions:
    \begin{itemize}
      \item[(a)] Multi-actor interaction: state space $2^{MN}$ vs.\ $2^N$
      \item[(b)] Endogenous landscape co-evolution via inPop and cycle4
            network statistics
      \item[(c)] Bounded rationality: logit precision $\beta$ nests the
            spectrum from random ($\beta=0$) to greedy ($\beta\to\infty$)
    \end{itemize}

  \item \textbf{Theorem 3 (Approximation):} Three independent proof
    strategies:
    \begin{itemize}
      \item[3a.] Constructive existence via dummy-coded dyadic covariates
      \item[3b.] Actor heterogeneity substitution ($Q \times M \geq 2^{K+1}$)
      \item[3c.] Universal approximation via mixed logit (McFadden \& Train 2000)
    \end{itemize}
\end{enumerate}


% ============================================================
%  APPENDIX B: Computational Demonstration
% ============================================================

\newpage
\part*{Appendix B: Computational Demonstration}
\addcontentsline{toc}{part}{Appendix B: Computational Demonstration}

\noindent\textit{The complete R script is contained in the companion file}
\texttt{nk\_saomnk\_computational\_demonstration.R}\textit{. Below we
summarize the protocol and expected results.}

\section*{B.1 Protocol}

The computational demonstration proceeds in five parts:

\begin{enumerate}
  \item \textbf{Standalone NK implementation.} A clean implementation of
    the NK adaptive walk (Kauffman 1993; Levinthal 1997) using standard
    functions: \texttt{powerkey()}, \texttt{make\_nk\_interaction\_matrix()},
    \texttt{nkland\_standard()}, \texttt{nk\_fitness()}, and
    \texttt{nk\_adaptive\_walk()}.

  \item \textbf{SaoMNK with $M=1$.} The \texttt{calc\_fit()} function
    extracted directly from the SaoMNK R6 class (line 4593 of
    \texttt{SAOM\_NK\_R6\_model.R}), plus a bridging function that constructs
    the equivalent $2^N \times N$ landscape matrix from standard NK payoff
    tables.

  \item \textbf{Equivalence tests.} For each $(N,K)$ pair in
    $\{(8,2), (8,4), (10,2), (10,5), (12,3), (12,6)\}$:
    \begin{itemize}
      \item 500 random landscapes $\times$ 200 adaptive walks each
      \item Point-wise fitness comparison on 10{,}000 random configurations
      \item Exhaustive local optima enumeration on 5 landscapes
      \item Two-sample Kolmogorov-Smirnov tests
    \end{itemize}

  \item \textbf{Publication-quality figures} (saved as PDF):
    \begin{itemize}
      \item \textbf{Figure 1:} NK vs.\ SaoMNK fitness scatter plots
            (perfect 45$^\circ$ line confirms exact equivalence)
      \item \textbf{Figure 2:} Paired density plots of final fitness with
            KS test $p$-value annotations
      \item \textbf{Figure 3:} Local optima counts across $(N,K)$ pairs
      \item \textbf{Figure 4:} ``Beyond NK'' --- $M \in \{1, 2, 5, 10\}$
            showing crowding, fitness variance, and trajectory volatility
    \end{itemize}

  \item \textbf{Summary statistics table} (LaTeX-formatted via
    \texttt{xtable}): mean/SD of steps to optimum, final fitness, KS
    statistics, and local optima counts.
\end{enumerate}

\section*{B.2 Expected Results}

Under the reduction conditions ($M=1$, $\boldsymbol{\theta}=\mathbf{0}$,
$\beta\to\infty$), the NK and SaoMNK implementations should produce:
\begin{itemize}
  \item \textbf{Identical fitness values} for all configurations (Figure 1:
    perfect correlation $r = 1.000$)
  \item \textbf{Indistinguishable walk statistics} (Figure 2: KS test
    $p \approx 1.0$)
  \item \textbf{Identical local optima} (Figure 3: zero difference)
\end{itemize}

For $M > 1$ (Figure 4), the ``Beyond NK'' results should show:
\begin{itemize}
  \item Increasing cross-actor fitness variance with $M$ (crowding effects)
  \item Non-monotonic fitness trajectories (landscape co-evolution)
  \item Greater trajectory volatility as actors' landscapes shift
    endogenously
\end{itemize}

\section*{B.3 Reproducibility}

All results use \texttt{set.seed(42)}. Required R packages:
\texttt{ggplot2}, \texttt{dplyr}, \texttt{tidyr}, \texttt{patchwork},
\texttt{xtable}, \texttt{scales}, \texttt{zoo}.


% ============================================================
%  APPENDIX C: Response to Anticipated Concerns
% ============================================================

\newpage
\part*{Appendix C: Response to Anticipated Concerns}
\addcontentsline{toc}{part}{Appendix C: Response to Anticipated Concerns}

\noindent\textit{The complete document is contained in the companion file}
\texttt{saomnk\_anticipated\_objections.tex}\textit{.
Below we list the eight concerns addressed.}

\bigskip

\begin{enumerate}
  \item \textbf{``The payoff structure is too structured / not random enough.''}
    --- The NK payoff table uses identical $U(0,1)$ random draws; SaoMNK
    generalizes the search process, not the landscape.

  \item \textbf{``This is just $M$ parallel NK searches.''}
    --- No. The inPop and cycle4 statistics create cross-actor dependencies
    producing endogenous landscape co-evolution absent in parallel NK.

  \item \textbf{``The logit choice model has the IIA problem.''}
    --- IIA is irrelevant under $\beta\to\infty$; for finite $\beta$, mixed
    logit (McFadden \& Train 2000) eliminates IIA while preserving
    tractability.

  \item \textbf{``The dynamics may not converge.''}
    --- $M=1$ converges by NK construction; $M>1$ has a Gibbs stationary
    distribution under Blume (1993) conditions; ergodicity on finite state
    space guarantees existence regardless.

  \item \textbf{``NK needs $N \cdot 2^{K+1}$ parameters. How does SAOM
    match this?''}
    --- Three mechanisms: constructive dummy covariates (exact but
    exponential), actor heterogeneity ($Q \times M \geq 2^{K+1}$), and
    universal approximation via mixed logit.

  \item \textbf{``This loses NK's analytical tractability.''}
    --- NK has no analytical tractability beyond $N=2, K=1$; all published
    results are simulation-based. SaoMNK ADDS estimability from empirical
    data, which NK lacks entirely.

  \item \textbf{``How is this different from ABMs that extend NK?''}
    --- SaoMNK is a proper statistical model (exponential family), not an
    ad hoc simulation. Parameters are estimable, interpretable, and
    testable. NK is a proven mathematical subset, not an analogy.

  \item \textbf{``The bipartite formulation changes NK fundamentally.''}
    --- The bipartite matrix reveals structure already implicit in NK.
    When $M=1$, the single row $\mathbf{b}_1$ IS the NK bit-string.
    The formulation adds competitive/cooperative dynamics central to
    strategy research.
\end{enumerate}


% ============================================================
%  FILE MANIFEST
% ============================================================

\newpage
\part*{File Manifest}
\addcontentsline{toc}{part}{File Manifest}

\begin{table}[h!]
\centering
\begin{tabular}{@{}lp{8cm}@{}}
\toprule
\textbf{File} & \textbf{Description} \\
\midrule
\texttt{saomnk\_online\_appendix\_master.tex} & This document (master appendix) \\
\texttt{saomnk\_nk\_equivalence\_proof.tex} & Appendix A: Formal proofs (Theorems 1--3) \\
\texttt{nk\_saomnk\_computational\_demonstration.R} & Appendix B: R script for computational verification \\
\texttt{saomnk\_anticipated\_objections.tex} & Appendix C: Response to anticipated concerns \\
\texttt{figures/} & Directory for computational demonstration output (Figures 1--4) \\
\bottomrule
\end{tabular}
\caption{Files in the \texttt{proofs/} directory.}
\end{table}


\end{document}
